% Owned proposed fragment B.135. Central integration is root-owned.
\subsection*{The complete marked inner star and an exact filled face}

\statement{B.135}{Proved}{The actual whole group has a marked inner-star relative presentation whose base map is noninjective in both categories; on an explicit retract block its kernel and the new source fixed subgroup are exact}
\label{stmt:B135}
Use the actual one-neck marking and source operators of B.123.
In the minus local tuple write
\[
 E_0=C_2(C)*F(a,b,d,e,f)*C_2(\tau),\qquad
 \delta=\operatorname{Inn}(\tau),\qquad a^*=\tau a\tau.
\]
The original orientation sends both cones to one in $C_2$ and
all five ordinary letters, and both source letters $m,v$, to zero.
Let $\Gamma_0$ be the actual ordinary group, and put
$E_{\rm act}=\Gamma_0\rtimes\langle\tau\rangle$.
All tuples, node words and operators used below are specified by
finite word formulas in the proof.

Let $R_-$ be the actual minus annular vertex, including its ten
node relations and their full $m$-orbit. Before the remaining
plus filling, the actual marked presentation is the proper amalgam
\begin{equation}\label{eq:B135-inner-star}
 G_{\rm in}=R_-*_{\langle m\rangle}
                  (\langle m\rangle\times C_2(\tau)).
\end{equation}
Its selected-cone centralizers, discretely and after full completion, are
\[
 C_{G_{\rm in}}(\tau)=\langle m\rangle\times\langle\tau\rangle,
 \qquad
 C_{\widehat G_{\rm in}}(\tau)
      =\overline{\langle m\rangle}\times\langle\tau\rangle.
\]
Thus the entire original ordinary inner subgroup has fixed subgroup
$\langle m\rangle$, and its actual full completion has fixed subgroup
$\overline{\langle m\rangle}\cong\widehat{\mathbb Z}$.
If $\beta$ is the actual $v$-operator and $\mathcal R_+$ the ten
actual plus orbifold node words, the \emph{whole} actual group is
\begin{equation}\label{eq:B135-whole-relative}
 E_{\rm act}=
 \left\langle G_{\rm in},v\ \middle|\
 v yv^{-1}=\beta(y)\ (y=a,b,d,e,f,C,\tau),\quad
 \mathcal R_+\right\rangle .
\end{equation}
In particular this presentation retains all signed source translates
of both node families, as well as every fiber conjugate.
It does not assert that $\beta$ descends to an automorphism of
$G_{\rm in}$.

The element
\begin{equation}\label{eq:B135-killed-word}
 w=[a,a^*]=a\tau a\tau a^{-1}\tau a^{-1}\tau
\end{equation}
is nontrivial in $G_{\rm in}$ and in its full completion, but
maps to one in $E_{\rm act}$ and its full completion. These
noninjectivity assertions also hold on the original ordinary
index-two subgroup and its full completion.

The following sharper comparison concerns actual subgroup blocks.
There are retractions onto
\begin{equation}\label{eq:B135-blocks}
 \begin{gathered}
 B=\langle a,\tau,m\rangle\leq G_{\rm in},\qquad
 B\cong(\mathbb Z(a)*C_2(\tau))\times\mathbb Z(m),\\
 Q=\langle a,\tau,m\rangle\leq E_{\rm act},\qquad
 Q\cong(\mathbb Z^2(a,a^*)\rtimes_{\rm swap}C_2(\tau))
                  \times\mathbb Z(m).
 \end{gathered}
\end{equation}
Writing $B_0=B\cap G_{\rm in}^0$ and $Q_0=Q\cap\Gamma_0$,
\begin{equation}\label{eq:B135-block-kernels}
 \begin{gathered}
 B_0=F(a,a^*)\times\mathbb Z(m),\qquad Q_0=\mathbb Z^3(a,a^*,m),\\
 \ker(B\to E_{\rm act})=\langle\!\langle w\rangle\!\rangle_B,
 \qquad
 \ker(B_0\to\Gamma_0)=\langle\!\langle w\rangle\!\rangle_{B_0}.
 \end{gathered}
\end{equation}
Each displayed subgroup has its full induced profinite topology
in the indicated ambient group. The corresponding kernels on
their actual completed closures are precisely the closed normal
closures of $w$.

The retraction $\Gamma_0\to Q_0$ is $\delta$-equivariant.
For the corrected actual contact
$A_s=s_1s_4^{-1}s_1^{-1}$ of B.130, the fixed subgroups on this
block are exactly
\begin{equation}\label{eq:B135-fixed-creation}
 \begin{gathered}
 B_0^\delta=\langle m\rangle,\qquad
 (\widehat B_0)^\delta=\overline{\langle m\rangle},\\
 Q_0^\delta=\langle aa^*,m\rangle
             =\langle A_s,m\rangle\cong\mathbb Z^2,\qquad
 (\widehat Q_0)^\delta
             =\overline{\langle A_s,m\rangle}\cong\widehat{\mathbb Z}^{\,2}.
 \end{gathered}
\end{equation}
Thus the filling creates this known source fixed norm. These
statements do not determine the whole kernel of
$G_{\rm in}\to E_{\rm act}$ or the full fixed subgroup of
$\Gamma_0$ or $\widehat\Gamma_0$.

\begin{proof}
\emph{B.135.1. Complete finite marking of both operators and both node lists.}
Recover the common neck coordinates from the displayed minus tuple by
\[
 B=CbC,\qquad c=aCa^{-1},\qquad K=F(a,B,d,e,f).
\]
The capital letter $B$ in this formula is a neck generator; the subgroup
$B$ in \eqref{eq:B135-blocks} will always be indicated as a subgroup.
The physical basis and actual mirror are
\begin{equation}\label{eq:B135-mirror}
 \begin{gathered}
 (O_1,O_2,O_3,O_4,O_5)
 =(aBa^{-1},a,fde d^{-1}f^{-1},d,d^{-1}fd),\\
 P_i=O_1\cdots O_i,\quad P_0=1,\qquad
 \eta(O_i)=P_{i-1}O_i^{-1}P_{i-1}^{-1}.
 \end{gathered}
\end{equation}
Its inverse is itself. The inverse physical basis change is
\[
 a=O_2,\quad B=O_2^{-1}O_1O_2,\quad d=O_4,\quad
 f=O_4O_5O_4^{-1},\quad e=O_5^{-1}O_4^{-1}O_3O_4O_5.
\]
Define the entire plus tuple by
\begin{equation}\label{eq:B135-plus-tuple}
 \begin{gathered}
 (a_+,B_+,d_+,e_+,f_+)
       =(\eta(a),\eta(B),\eta(d),\eta(e),\eta(f)),\\
 C_+=a_+^{-1}\tau a_+,\qquad b_+=C_+B_+C_+.
 \end{gathered}
\end{equation}
The ordered systems $(a,b,d,e,f,C,\tau)$ and
$(a_+,b_+,d_+,e_+,f_+,C_+,c)$ are complete free-product bases
of $E_0$. To invert the latter, recover $B_+=C_+b_+C_+$,
apply $\eta$ to its five ordinary coordinates, and use
$C=a^{-1}ca$ and $\tau=a_+C_+a_+^{-1}$.
The other cone $c$ remains part of
that basis. Thus the following formulas specify automorphisms
of the whole group, rather than actions on omitted peripherals.

In the minus basis the actual $m$-action is
\begin{equation}\label{eq:B135-m-action}
 \alpha_m(C)=f^{-1}Cf,\quad
 \alpha_m(f)=f^{-1}CfCf,\quad
 \alpha_m(U)=U\ (U=a,b,d,e),\quad \alpha_m(\tau)=\tau.
\end{equation}
It is the partial inner automorphism $\operatorname{Inn}(f^{-1}C)$
on $\langle C,f\rangle$, with inverse partial
$\operatorname{Inn}(Cf)$ and all outsiders fixed.
Define $\alpha_p$ by the same two formulas on
$(C_+,f_+)$, fixing $a_+,b_+,d_+,e_+,c$.
Then the full actual source correction is
\begin{equation}\label{eq:B135-beta}
 q_v=a_+f_+a_+^{-1},\qquad
 \beta=\operatorname{Inn}(q_v)\alpha_p,\qquad
 \beta^{-1}=\alpha_p^{-1}\operatorname{Inn}(q_v^{-1}).
\end{equation}
These are the actual $v=q_vp$ and inverse operators of B.123.
In particular $\beta(\tau)=\tau$; both signed operators
commute with $\delta$ on the ordinary subgroup.

For any local tuple $(A,b_0,D,E,f_0,C_0)$ let
\begin{equation}\label{eq:B135-ten-words}
 \begin{gathered}
 \mathcal R(A,b_0,D,E,f_0,C_0)=
 \bigl\{[U,C_0UC_0]:U=A,b_0,D,E\bigr\}\\
 {}\cup\bigl\{[A,b_0],[C_0AC_0,D],[C_0AC_0,E],
 [D,C_0EC_0],\ [b_0,f_0Df_0^{-1}],\ [b_0,f_0Ef_0^{-1}]\bigr\}.
 \end{gathered}
\end{equation}
Take $\mathcal R_- =\mathcal R(a,b,d,e,f,C)$ and
$\mathcal R_+=\mathcal R(a_+,b_+,d_+,e_+,f_+,C_+)$.
These are the ten actual based words in each region from B.123,
not only their peripheral conjugacy classes. The local replacement
of $hdh^{-1}$ by $fdf^{-1}$, and similarly for $e$, uses the
actual generic equations $h=fm$ and $[m,d]=[m,e]=1$.
The two ordinary sheets of every one of these twenty words
give all forty ordinary nodes. No sheet is discarded here.

\emph{B.135.2. The exact generic star and the complete relative presentation.}
Set $u=Cf$ and introduce the auxiliary letter $k=um$.
This $k$ is \emph{odd} in the original orientation and is not
the even far/base generator called $k$ in B.112; here that
even generator is $m$.
Equations \eqref{eq:B135-m-action} give $\alpha_m(u)=u$ and
$\alpha_m(C)=u^{-1}Cu$. Consequently the generic $m$ mapping
torus has the inverse Tietze presentation
\begin{equation}\label{eq:B135-graph-star}
 \left\langle C,\tau,m,k,a,b,d,e\ \middle|\
 C^2=\tau^2=1, [m,k]=[k,C]=[m,\tau]=1,
 [m,U]=1\ (U=a,b,d,e)\right\rangle,
\end{equation}
with inverse substitution $f=Ckm^{-1}$.
Conversely the two commutations involving $k$ recover
$\alpha_m(u)=u$ and $\alpha_m(C)=u^{-1}Cu$, hence both
images of $C,f$ in \eqref{eq:B135-m-action}.

Every minus node lies in $C_2(C)*F(a,b,d,e,f)$ and contains
no $\tau$. Imposing its ten words in the first factor, with
their full $m$-orbit supplied by conjugation, therefore gives
exactly \eqref{eq:B135-inner-star}, by the two inverse
presentations. The common $\langle m\rangle$ embeds: the
actual base-winding homomorphism on $R_-$ sends $m$ to one
in $\mathbb Z$ and every fiber generator to zero.
This proves properness in the discrete category.

In its Bass--Serre tree, $\tau$ fixes the second vertex but
no incident edge. Indeed an edge stabilizer is conjugate to
the torsion-free group $\langle m\rangle$, whereas $\tau$
has order two. Its fixed subtree is therefore a single vertex.
Every centralizer element preserves that vertex, proving the
centralizer formula in the statement.

Start now from the actual pre-node presentation
$E_0\rtimes F(m,v)$ of B.123 and impose both lists
\eqref{eq:B135-ten-words}. Grouping its $m$-equations and
minus nodes first gives $G_{\rm in}$. All remaining equations
are precisely the seven $v$-equations and the ten plus nodes
in \eqref{eq:B135-whole-relative}. Expanding the presentation
of $G_{\rm in}$ reverses this grouping and fixes all marked
generators. This proves the exact whole relative presentation.
Normal closure in this presentation includes conjugation by
every signed word in $m,v$ and by every fiber word. There is
no length bound, finite orbit substitution, or assertion that
the remaining equations preserve the inner base injectively.

\emph{B.135.3. A simultaneous ordinary basis and the actual characters.}
In the ordinary free group of rank eleven put
\begin{equation}\label{eq:B135-new-basis}
 \begin{gathered}
 z=\tau C,\quad p_0=\tau u,\quad q=u\tau,\quad
 r=zq=\tau f\tau,\quad s=p_0q,\\
 \mathcal Y=(r,s,q,a,b,d,e,a^*,b^*,d^*,e^*),\qquad U^*=\tau U\tau.
 \end{gathered}
\end{equation}
The standard ordinary Schreier basis is
$(z,f,\tau f\tau,U,U^*)$. The inverse changes
\[
 z=rq^{-1},\qquad p_0=sq^{-1},\qquad
 f=z^{-1}p_0=q r^{-1}s q^{-1}
\]
show that \(\mathcal Y\) is a simultaneous free basis.
Substitution in \eqref{eq:B135-m-action} and conjugation by
$\tau$ give the full formulas
\begin{equation}\label{eq:B135-basis-actions}
 \begin{gathered}
 \alpha_m(r)=r^{-1}s,\quad
 \alpha_m(s)=s,\quad\alpha_m(q)=q,\quad
 \alpha_m(U)=U,\quad\alpha_m(U^*)=U^*,\\
 \delta(r)=q r^{-1}s q^{-1},\quad
 \delta(s)=q s q^{-1},\quad\delta(q)=s q^{-1},\quad
 \delta(U)=U^*,\quad\delta(U^*)=U.
 \end{gathered}
\end{equation}
The inverse $m$-action sends $r$ to $sr^{-1}$ and fixes
the other ten letters. For every integer $n$,
\[
 \alpha_m^{2n}(r)=s^{-n}rs^n,\qquad
 \alpha_m^{2n+1}(r)=s^{-n}r^{-1}s^{n+1}.
\]
These identities follow from $\alpha_m^2$ and its inverse,
so include negative indices without an observed-orbit argument.

The abelian action of \emph{each} of $\alpha_m$ and $\beta$
on this basis is
\begin{equation}\label{eq:B135-abelian-actions}
 [r]\longmapsto-[r]+[s],\qquad
 [y]\longmapsto[y]\quad(y\in\mathcal Y\setminus\{r\}).
\end{equation}
For $\beta$ this is the direct exponent calculation in the
finite word rules \eqref{eq:B135-mirror}--\eqref{eq:B135-beta};
the complete actual-word certificate checks all eleven columns,
both signed inverse operators and their ordinary/orbifold
diagrams. Thus the $a$ and $a^*$ exponent characters are
invariant under every signed source word. All forty nodes
are commutators; their entire source orbit and normal closure
have zero exponent in every ordinary fiber coordinate.
Both characters therefore descend through the whole actual
filling. Together with the $m$-exponent in the actual free
source base they give
\begin{equation}\label{eq:B135-rho}
 \rho:\Gamma_0\longrightarrow\mathbb Z^3,\qquad
 \rho=(\operatorname{exp}_a,\operatorname{exp}_{a^*},
                         \operatorname{exp}_m).
\end{equation}
The computation also gives $\Gamma_0^{\rm ab}=\mathbb Z^{12}$:
the sole fiber-coinvariant relation is $[s]=2[r]$, followed
by the two independent source-base classes. We will use the
explicit section below, rather than abelianization alone,
to prove an embedded torus group.

\emph{B.135.4. Actual plus nodes fill one specified face.}
For clarity retain the entire based conjugation in the
four-cross argument of B.130. Set
\[
 A=a_+,\quad C_0=C_+,\quad b_0=b_+,\quad
 V=AB_+A^{-1},\quad r_0=[A,C_0AC_0].
\]
The physical prefix rule gives the literal words
\[
 V=aB^{-1}a^{-1},\qquad A=V^{-1}a^{-1}V,\qquad
 VA^{-1}V^{-1}=a.
\]
Three entries of the actual plus list are
$r_0,[b_0,C_0b_0C_0],[A,b_0]$. Free reduction using cone
squares gives
\begin{equation}\label{eq:B135-four-cross}
 \begin{aligned}
 Ar_0A^{-1}&=[A,\delta(A)],\\
 \bigl(A[b_0,C_0b_0C_0]A^{-1}\bigr)^{-1}
      &=[V,\delta(V)],\\
 A[A,b_0]A^{-1}&=[A,\delta(V)],\\
 \delta\bigl(A[A,b_0]A^{-1}\bigr)^{-1}
      &=[V,\delta(A)].
 \end{aligned}
\end{equation}
Hence $D=\langle V,A\rangle$ commutes elementwise with
$\delta(D)$ in the actual final group. Since $a\in D$,
we have $[a,a^*]=1$. This implication is inside the whole
defining presentation and assumes no surviving regional
embedding. The norm $N_\delta(t)=t\delta(t)$ is a
homomorphism on $D$. The literal raw source forms are
\[
 s_1=A(C_0b_0C_0b_0)A^{-1},\qquad
 s_4=A(C_0b_0r_0b_0^{-1}C_0\,AC_0AC_0)A^{-1}.
\]
They satisfy $s_1=N_\delta(V)$ and
$s_4=(AC_0b_0)r_0(AC_0b_0)^{-1}N_\delta(A)$ before
the node is killed. Thus after the actual node relation
\[
 A_s=s_1s_4^{-1}s_1^{-1}
    =N_\delta(VA^{-1}V^{-1})=aa^*.
\]
The literal $m$-action fixes both $a$ and $a^*$ by
\eqref{eq:B135-basis-actions}. Consequently
\[
 i:\mathbb Z^3\longrightarrow\Gamma_0,\qquad
 i(j,k,\ell)=a^j(a^*)^k m^\ell
\]
is a homomorphism, and \eqref{eq:B135-rho} gives $\rho i=1$.
This proves the actual $\mathbb Z^3$ embedding and retraction.
The involution swaps the first two coordinates and fixes the
third, both on the domain and the codomain. Thus the retraction
is $\delta$-equivariant. The four corrected contacts have
images $(1,1,0),0,0,0$, while $v$ maps to zero and $m$ to
$(0,0,1)$. The actual source image maps onto the entire fixed
part of this retracted face.

\emph{B.135.5. Nonzero kernel before and after full completion.}
The element $a\in R_-$ has winding zero and is nontrivial:
its exponent character is preserved by the generic action
and vanishes on every minus commutator. Therefore
$a\notin\langle m\rangle$. Also $\tau$ is outside that
edge. The word \eqref{eq:B135-killed-word} has reduced
alternating amalgam length eight in \eqref{eq:B135-inner-star},
so it is nontrivial there, and its original orientation is even.
Section B.135.4 proves that its actual final image is one.

For a direct full-completion witness send
\[
 a\longmapsto(12),\qquad\tau\longmapsto(23),\qquad
 b,d,e,f,C,m\longmapsto1
\]
to $S_3$. On $R_-$ this is its $a$-exponent modulo two;
all generic equations and all ten minus nodes hold. The
common $m$ has the same image on the second vertex, and
$\tau^2=[m,\tau]=1$ hold. Thus this is a homomorphism of
the \emph{whole} inner star. Its image of $w$ is a nontrivial
three-cycle. It extends to $\widehat G_{\rm in}$ and restricts
to a finite witness on the original ordinary subgroup.
This proves the stated completed noninjectivity without a
profinite-amalgam normal-form assumption.

There is also an exact full-inner compact centralizer calculation.
Let finite targets of $R_-$ and $\langle m\rangle\times C_2(\tau)$
be prescribed. Choose a positive integer $N$ divisible by the
orders of both prescribed images of $m$. For the first vertex
combine its target with the actual winding modulo $N$; for the
second combine its target with its $m$-coordinate modulo $N$.
Their restrictions to the edge have the identical exact kernel
$N\mathbb Z$, and the original targets are retained. The finite
vertex images thus contain the same embedded cyclic edge group
$C_N$.

Take suitable multiples of the two finite regular permutation
actions on a common finite set. The restrictions to $C_N$ are
free actions with equal orbit counts; identifying corresponding
orbits equivariantly makes them agree. This gives a finite
quotient of the amalgam which is injective on both finite vertex
images. It follows that both vertices and the edge have their
full induced topology, and that the universal full completion is
the proper profinite amalgam
\begin{equation}\label{eq:B135-compact-star}
 \widehat G_{\rm in}
   =\widehat R_-\amalg_{\widehat{\mathbb Z}(m)}
          \bigl(\widehat{\mathbb Z}(m)\times C_2(\tau)\bigr).
\end{equation}
This proves properness from actual finite targets; it does not
presume that an arbitrary completed amalgam is proper.

Use its standard completed coset tree
\citep[Proposition~3.8]{B_ZalesskiiMelnikov1989}.
The involution $\tau$ fixes the second vertex and no edge:
every edge stabilizer is a conjugate of the torsion-free
$\widehat{\mathbb Z}(m)$. Its nonempty fixed subgraph is connected
by the prime-order fixed-tree theorem
\citep[Theorem~2.8]{B_ZalesskiiMelnikov1989}.
A connected profinite graph without edges has one vertex.
Every compact centralizer element therefore preserves that vertex,
so its full centralizer is exactly
$\widehat{\mathbb Z}(m)\times C_2(\tau)$.
The original ordinary subgroup has finite index, hence its full
completion is exactly the orientation kernel in
\eqref{eq:B135-compact-star}. Intersecting the centralizer with
this kernel gives precisely $\widehat{\mathbb Z}(m)$.

The image of $aa^*$ is the nontrivial actual fixed source element
$A_s$, in both categories, by the actual retracted face. It has
zero base winding. If it were the image of an inner fixed element
$m^j$ or $m^\lambda$, respectively, the actual base character
would force $j=0$ or $\lambda=0$. Its nontriviality rules this out.
Thus $A_s$ has no fixed preimage from the \emph{entire} ordinary
inner group or its full completion. This is a precise failure of
fixed lifting through this filling, for a known source element;
it is not a fixed element outside the final actual source.

\emph{B.135.6. The exact block kernels and their full topology.}
Kill $C,f,b,d,e$ in the complete inner presentation. All
generic equations survive, and every minus node becomes one;
the only possibly surviving diagonal is $[a,CaC]$, which
becomes $[a,a]$. This gives the actual retraction onto
\[
 B=(\mathbb Z(a)*C_2(\tau))\times\mathbb Z(m)
\]
with the evident generator section. The original orientation
restricts to $\tau\mapsto1$, $a,m\mapsto0$, so the actual
ordinary subgroup is $B_0=F(a,a^*)\times\mathbb Z(m)$.
The retraction from all $G_{\rm in}$ need not preserve the
\emph{original} orientation. Full topology of $B_0$ nonetheless
follows from full topology of the actual retract $B$ and
finite index in $B$, then restriction to the ordinary ambient
subgroup. No orientation-preserving retraction is assumed.

The actual equality $w=1$ gives a map from
\[
 Q=B/\langle\!\langle w\rangle\!\rangle_B
   =(\mathbb Z^2(a,a^*)\rtimes C_2(\tau))\times\mathbb Z(m)
\]
into $E_{\rm act}$. Extend the equivariant retraction
\eqref{eq:B135-rho} to $E_{\rm act}=\Gamma_0\rtimes C_2(\tau)$
by sending $\tau$ to the displayed swap involution. Its
composite with this map is the identity on $a,\tau,m$.
Thus $Q$ embeds as precisely the actual image of $B$, with
an actual retraction; its ordinary subgroup is $Q_0=\mathbb Z^3$.
This proves both discrete kernel formulas
\eqref{eq:B135-block-kernels}. For the ordinary one, the
presentation $F(a,a^*)\times\mathbb Z(m)\to\mathbb Z^3$
has exactly the normal commutator relation $w$.

Retractions induce injections on full completions, with image
the actual closures, and finite-index subgroups inherit full
topology by taking finite-index cores. Hence the completed
map $\widehat B\to\widehat E_{\rm act}$ factors through
$\widehat Q$ followed by an injection, and similarly for
$B_0,Q_0$. Finite quotients of $B/\langle\!\langle w\rangle\!\rangle$
are exactly the finite quotients of $B$ killing $w$.
The universal property therefore gives
\begin{equation}\label{eq:B135-closed-kernels}
 \begin{gathered}
 \ker(\widehat B\to\widehat E_{\rm act})
   =\overline{\langle\!\langle w\rangle\!\rangle}^{\,\widehat B},\\
 \ker(\widehat B_0\to\widehat\Gamma_0)
   =\overline{\langle\!\langle w\rangle\!\rangle}^{\,\widehat B_0}.
 \end{gathered}
\end{equation}
The bars denote closed normal closures in the stated actual
embedded completions. No completion of an unspecified
infinite kernel has been substituted.

\emph{B.135.7. The exact fixed change on the block and its scope.}
The swap action on $F(a,a^*)$ has trivial fixed subgroup.
Indeed its extension by the swap involution is
$C_2(\tau)*\mathbb Z(a)$; the factor-centralizer normal form
gives $C(\tau)=\langle\tau\rangle$, with trivial intersection
with the orientation kernel. In the full profinite category
the corresponding universal presentation is
$C_2(\tau)\amalg\widehat{\mathbb Z}(a)$.
The same centralizer equality is Theorem B' of
\citet{A_HR1985}, in precisely the full-profinite scope used
in B.129. This proves the two pre-filling fixed formulas in
\eqref{eq:B135-fixed-creation}. After filling, the involution
on $\mathbb Z^3$ or $\widehat{\mathbb Z}^{\,3}$ swaps the
first two coordinates and fixes the third, proving the two
post-filling formulas directly.

For every signed source word $t\in F(m,v)$ there is also the
exact generic identity
\[
 \alpha_t([a,\delta(a)])
   =[\alpha_t(a),\delta(\alpha_t(a))].
\]
Its final norm is
$\alpha_t(aa^*)=\alpha_t(A_s)$, an actual source conjugate.
This retains the whole source orbit and does not identify
all nodes with that one orbit. The full inner-star base map
is already \emph{noninjective}, by B.135.5. What remains
undetermined is its whole kernel and the resulting global
fixed-source comparison. The proved whole-inner compact
centralizer supplies the specific negative lifting statement
of B.135.5; it gives no positive fixed lifting or centralizer
formula for the final ambient group.
\end{proof}

\paragraph{Complete actual verification and preserved records.}
The literal basis producer checks all twenty-two inverse
basis compositions, all forty-four signed operator inverse
comparisons, the involution square and forty-four exact
source/involution comparisons, all forty whole node words
with both peripheral factors and round trips, all eleven
orbifold/ordinary marking diagrams, and all six raw source
round trips. The relative-face companion checks all four
orbifold operators, all twenty full node words, both complete
abelian matrices, the involution and contact characters.
The kernel companion checks all nineteen inner defining
relations in the single specified $S_3$ target and its
nonidentity value on $w$. These are complete word and
presentation checks; none asserts closure of a finite orbit.
The first basis run's caught inverse-$B$ transcription error
is preserved with its genuine failing output; the corrected
inverse and every subsequent full comparison pass.
\begingroup\raggedright\footnotesize
\path{work_CE/code/compile_m_star_relative_presentation27.py},
\path{work_CE/code/compile_m_star_relative_presentation27.out},
\path{work_CE/next_symmetric_nodes/m_star_relative_presentation27.json};
\path{work_CE/code/check_m_star_relative_face27.py},
\path{work_CE/code/check_m_star_relative_face27.out},
\path{work_CE/next_symmetric_nodes/m_star_relative_face27.json};
\path{work_CE/code/check_inner_star_kernel27.py},
\path{work_CE/code/check_inner_star_kernel27.out},
\path{work_CE/next_symmetric_nodes/inner_star_kernel27.json}.
Full research proofs:
\path{work_CE/whole_node_m_star_presentation27.txt},
\path{work_CE/inner_star_actual_kernel27.txt},
\path{work_CE/inner_star_face_exact_kernel27.txt}.
\par\endgroup
